\usepackage[top=2.5cm, bottom=3cm, left=3cm, right=3cm]{geometry}
\usepackage[spanish, main=english]{babel}
\usepackage[pdftex]{graphicx}
\usepackage{hyperref}
   \hypersetup{colorlinks, linkcolor=blue, urlcolor=blue, citecolor={red}}
\usepackage{color}
\usepackage[capitalize]{cleveref}

% --- Configuración de Encabezados Dinámicos ---
\usepackage{fancyhdr}

\definecolor{verdeUdeA}{rgb}{0.325,0.506,0.208} 

\pagestyle{fancy}
\fancyhf{} 

\renewcommand{\chaptermark}[1]{\markboth{#1}{}} 

% Lógica mejorada:
% Si el contador de capítulos es > 0, muestra el nombre del capítulo (\leftmark)
% Si es 0 (estamos en el resumen, TOC, etc.), muestra el título corto (\smalltitulo)
\lhead{\footnotesize\MakeUppercase{%
  \ifnum\value{chapter}>0
    \leftmark%
  \else
    \smalltitulo%
  \fi
}}

\chead{}
\rhead{\thepage} 

\renewcommand{\headrule}{{\color{verdeUdeA}\rule{\textwidth}{1pt}}}
\renewcommand{\headrulewidth}{2pt}
\cfoot{}

% % --- Configuración de Encabezados Dinámicos e Inteligentes ---
% \usepackage{fancyhdr}
% \usepackage{etoolbox} % Necesario para la lógica IF

% \definecolor{verdeUdeA}{rgb}{0.325,0.506,0.208} 

% \pagestyle{fancy}
% \fancyhf{} % Limpia encabezados y pies anteriores

% % 1. Definimos cómo se guarda la marca del capítulo
% % Esto hace que \leftmark contenga solo el título del capítulo
% \renewcommand{\chaptermark}[1]{\markboth{#1}{}} 

% % 2. Lógica del Encabezado Izquierdo (\lhead)
% % IF: Si \leftmark está vacío (antes del primer capítulo), muestra \smalltitulo
% % ELSE: Si ya hay un capítulo, muestra \leftmark
% \lhead{\footnotesize\MakeUppercase{%
%   \expandafter\ifstrempty\expandafter{\leftmark}{%
%     \smalltitulo%
%   }{%
%     \leftmark%
%   }%
% }}

% \chead{}
% \rhead{\thepage} % Número de página a la derecha

% % 3. Línea verde de la UdeA
% \renewcommand{\headrule}{{\color{verdeUdeA}\rule{\textwidth}{1pt}}}
% \renewcommand{\headrulewidth}{2pt}
% \cfoot{}